

It is sometimes necessary to control the finite field 
that is used in the construction of a matrix group.
For prime fields, this is not an issue. 
For extension fields, the choice of polynomial does matter, as 
the generators depend on specific choices made for the finite field. 
Magma and GAP use Conway polynomials, which are difficult to compute.
Orbiter has a built-in table of primitive polynomials. 
As explained in Section~\ref{sec:extension:fields},
Orbiter allows to specify the polynomial that should be used to create the finite field.
The next example shows an instance where choosing 
the polynomial is important. 
We are recreating a group 
from the electronic Atlas on finite simple groups~\cite{Atlas3}.

\bigskip

The electronic Atlas of finite simple groups~\cite{Atlas3} 
lists generators for $U_3(3)$ as $3 \times 3$ matrices over the field $\bbF_9$
using the following short Magma~\cite{Magma} program:

\orbiteroutput{GRAPHICS/LATEX/U_3_3_magma}

The generators are given using the Magma command \verb'CambridgeMatrix', 
which allows for more efficient coding of field elements.
The field elements are coded as base-3 integers (like in Orbiter) 
with respect to the Magma version of $\bbF_9.$ 
The polynomial for $\bbF_9$ can be determined using the 
following Magma command: 
%which can be typed into Magma (or 
%the free Magma online calculator at~\cite{MagmaOnline}):

\orbiteroutput{GRAPHICS/LATEX/magma_poly_F9}

It results in
\orbiteroutput{GRAPHICS/LATEX/magma_poly_F9_out}
which is the Magma way of printing the polynomial $X^2+2X+2.$ 
If $\alpha$ is a root of the polynomial over $\bbF_3$, then
$$
\alpha^2 = \alpha+1.
$$
The coefficient vector of 
the polynomial is $(1,2,2).$ 
As an integer written in base-3, we 
obtain 
$$
1\cdot 3^2 + 2 \cdot 3 + 2 = 17.
$$
The desired subgroup can now be created using the command

\sourcecode{U_3_3}

Group theoretic activities 
will be discussed in Section~\ref{sec:group:theoretic:activities}.

\bigskip

As an example of a large group, consider the Conway group ${\rm Co}_3.$ 
Following~\cite{SuleimanWilson1997}, the group can be generated using two 
matrices of dimension $22$ over $\bbF_2$. 
We use the makefile variables to give each generator in compact form. 

\sourcecodevariable{CONWAY_GEN1}


\sourcecodevariable{CONWAY_GEN2}

Then we define vectors for each of the generators. 
We concatenate the two generators to 
form one long vector, 
which is passed to the \verb'-subgroup_by_generators' command.
Finally, we create a report for the group.



\sourcecode{Co3}



\bigskip

The next example creates the Ree group 
in 7 dimensions over the field $\bbF_{27}$. 
Again, we use makefile variables to specify the 
two generators as $7 \times 7$ matrices over $\bbF_{27}$ 
and concatenate them, 
before passing them to the \verb'-subgroup_by_generators' command.

\sourcecodevariable{Ree_gen1}


\sourcecodevariable{Ree_gen2}


\sourcecode{Ree_27}


\bigskip


Let us take a look at the Higman-Sims group. 
It is a permutation group acting on 100 points. 
More precisely, the group is related 
to the automorphism group of a certain graph that arises from 
the Witt design on 22 points. 
For a construction of the design, see Section~\ref{sec:design:assuming:symmetry}.
For a construction of the graph, see Section~\ref{sec:graphtheory}.

\bigskip


We use the following three makefile variables 
(for the sake of space, the third one, \verb'HIGMAN_SIMS_GENS' is not shown), 

%\sourcecodevariable{HIGMAN_SIMS_GENS}

\sourcecodevariable{HIGMAN_SIMS_NB_GENS}

\sourcecodevariable{HIGMAN_SIMS_GO}


The following command creates the group in the permutation representation:



\sourcecode{Higman_Sims_group}

\bigskip


The next command computes the conjugacy classes of elements:

\sourcecode{Higman_Sims_group_classes}

\bigskip

We also compute the conjugacy classes of subgroups:

\sourcecode{Higman_Sims_group_subgroup_lattice}

\bigskip

The next two commands define the representatives 
for the two distinct conjugacy classes of subgroups of order 252000:

\sourcecode{Higman_Sims_subgroup_176a}

\sourcecode{Higman_Sims_subgroup_176b}


The command to define a subgroup from 
the previously computed lattice of subgroups performs 
a shortening of the base. 
Base points whose basic orbit is trivial are dropped.


\bigskip



It is sometimes desirable to change the underlying form of a classical group. 
The following command creates generators for the group $\PSp(6,2)$ with respect to the form 
$$
\left[
\begin{array}{cc}
0 & I\\
I & 0\\
\end{array}
\right].
$$

\sourcecode{Butler12_PSP_6_2_conjugate}



\bigskip



The following command creates the subfield subgroup of the Weyl Group of type $E_6$.

\sourcecode{Hirschfeld_surface_get_subfield_subgroup}



\bigskip



